% Part: sets-functions-relations
% Chapter: relations
% Section: relations-as-sets

\documentclass[../../../include/open-logic-section]{subfiles}

\begin{document}

\olfileid{sfr}{rel}{set}
\olsection{সেট হিসেবে সম্পর্ক}

\begin{explain}
\olref[sfr][set][imp]{sec}-এ কয়েকটি গুরুত্বপূর্ণ সেটের কথা বলেছি:
$\Nat$, $\Int$, $\Rat$, $\Real$। এই সেটগুলির কোনও কোনওটির
!!{element}sের মধ্যে কিছু আকর্ষণীয় সম্পর্ক তোমার নিশ্চয়ই মনে আছে।
যেমন, প্রতিটি সেটেই একটি সুপরিচিত \emph{ক্রমসম্পর্ক} আছে।
এ ছাড়াও আছে \emph{অভিন্ন হওয়ার} সম্পর্ক: প্রতিটি বস্তু নিজের সঙ্গে
এই সম্পর্কে থাকে, অন্য কোনও বস্তুর সঙ্গে নয়। আরও অনেক আকর্ষণীয়
সম্পর্কের সঙ্গে আমাদের পরিচয় হবে; সম্ভাব্য সম্পর্কের সংখ্যা আরও বেশি।
তবে সেগুলির আলোচনা করার আগে দেখাব যে, সম্পর্ককে বিশেষ ধরনের
সেট হিসেবে দেখা যায়।

এর জন্য \olref[sfr][set][pai]{sec}-এর দুটি বিষয় মনে করো। প্রথমটি
\emph{ক্রমযুগলের} ধারণা: $a$ ও $b$ দেওয়া থাকলে $\tuple{a, b}$
গঠন করতে পারি। এখানে উপাদানের ক্রম \emph{অবশ্যই} গুরুত্বপূর্ণ।
তাই $a \neq b$ হলে $\tuple{a, b} \neq \tuple{b, a}$।
(ক্রমহীন যুগল, অর্থাৎ $2$-উপাদানবিশিষ্ট সেটের সঙ্গে তুলনা করো;
সেখানে $\{a, b\}=\{b, a\}$।) দ্বিতীয়টি \emph{কার্তেসীয় গুণফলের}
ধারণা: $A$ ও $B$ সেট হলে $A \times B$ গঠন করতে পারি।
এটি $x \in A$ এবং $y \in B$ এমন সমস্ত যুগল $\tuple{x, y}$-এর সেট।
বিশেষত, $A^{2}= A \times A$ হল $A$ থেকে গঠিত সমস্ত ক্রমযুগলের সেট।

এবার কোনও সেটের উপর একটি নির্দিষ্ট সম্পর্ক বিবেচনা করব:
স্বাভাবিক সংখ্যার সেট~$\Nat$-এর উপর $<$-সম্পর্ক। সংখ্যার সেই সমস্ত
যুগল $\tuple{n, m}$-এর সেট বিবেচনা করো, যাদের জন্য $n<m$। অর্থাৎ,
\[
  R=\Setabs{\tuple{n, m}}{n, m \in \Nat \text{ এবং } n<m}.
\]
$n$, $m$-এর চেয়ে ছোট হওয়া এবং $\tuple{n, m}$ যুগলটি $R$-এর সদস্য
হওয়ার মধ্যে ঘনিষ্ঠ যোগ আছে:
\[
      n<m\text{ হবে যদি এবং কেবল যদি }\tuple{n, m} \in R.
\]
বস্তুত, কোনও তথ্য না হারিয়েই $R$ সেটটিকে $\Nat$-এর উপর
$<$-সম্পর্ক \emph{হিসেবেই} বিবেচনা করতে পারি।

একইভাবে, সংখ্যার মধ্যে যে-কোনও সম্পর্কের জন্য $\Nat^{2}$-এর
একটি উপসেট গঠন করতে পারি। বিপরীতভাবে, সংখ্যার যুগলের যে-কোনও সেট
$S \subseteq \Nat^{2}$ দেওয়া থাকলে, সংখ্যার মধ্যে তার অনুরূপ একটি
সম্পর্ক থাকে: $n$, $m$-এর সঙ্গে সেই সম্পর্কে থাকবে যদি এবং কেবল যদি
$\tuple{n, m} \in S$ হয়। এর ভিত্তিতে নিচের সংজ্ঞা দিচ্ছি:
\end{explain}

\begin{defn}[দ্বিপদ সম্পর্ক]
কোনও সেট $A$-এর উপর \emph{দ্বিপদ সম্পর্ক} হল $A^{2}$-এর একটি উপসেট।
$R \subseteq A^{2}$ যদি $A$-এর উপর দ্বিপদ সম্পর্ক হয় এবং
$x, y \in A$ হয়, তবে $\tuple{x, y} \in R$ বোঝাতে কখনও
$Rxy$ (অথবা $xRy$) লিখি।
\end{defn}

\begin{ex}
  \ollabel{relations}
স্বাভাবিক সংখ্যার যুগলগুলির সেট $\Nat^{2}$-কে নিচের মতো
দ্বিমাত্রিক ম্যাট্রিক্সে সাজানো যায়:
\[
  \begin{array}{ccccc}
  \mathbf{\tuple{ 0,0 }} & \tuple{ 0,1 } &
    \tuple{ 0,2 } & \tuple{ 0,3 } & \ldots\\
  \tuple{ 1,0 } & \mathbf{\tuple{ 1,1 }} &
    \tuple{ 1,2 } & \tuple{ 1,3 } & \ldots\\
  \tuple{ 2,0 } & \tuple{ 2,1 } &
    \mathbf{\tuple{ 2,2 }} & \tuple{ 2,3 } & \ldots\\
  \tuple{ 3,0 } & \tuple{ 3,1 } & \tuple{ 3,2 } &
    \mathbf{\tuple{ 3,3 }} & \ldots\\
  \vdots & \vdots & \vdots & \vdots & \mathbf{\ddots}
  \end{array}
\]
এখানে কর্ণের যুগলগুলি মোটা হরফে দেখিয়েছি। কারণ, কর্ণের যুগলগুলি
নিয়ে গঠিত $\Nat^2$-এর উপসেটটি, অর্থাৎ
\[
  \{\tuple{0,0 }, \tuple{ 1,1 }, \tuple{ 2,2 }, \dots\},
  \]
হল \emph{$\Nat$-এর উপর অভিন্নতার সম্পর্ক}। (এই সম্পর্কটি বহুল
ব্যবহৃত বলে, যে-কোনও সেট $A$-এর জন্য
$\Id{A}=\Setabs{\tuple{ x,x }}{x \in A}$ সংজ্ঞায়িত করি।)
কর্ণের উপরের সমস্ত যুগলের উপসেটটি, অর্থাৎ
\[
  L = \{\tuple{ 0,1 },\tuple{ 0,2 },\ldots,\tuple{ 1,2 },
  \tuple{ 1,3 }, \dots, \tuple{ 2,3 }, \tuple{ 2,4 },\ldots\},
\]
হল \emph{চেয়ে ছোট হওয়ার} সম্পর্ক: $Lnm$ হবে যদি এবং কেবল যদি
$n<m$ হয়। কর্ণের নিচের যুগলগুলির উপসেটটি, অর্থাৎ
\[
  G=\{ \tuple{ 1,0 },\tuple{ 2,0 },\tuple{
    2,1 }, \tuple{ 3,0 },\tuple{ 3,1 },\tuple{ 3,2 }, \dots\},
\]
হল \emph{চেয়ে বড় হওয়ার} সম্পর্ক: $Gnm$ হবে যদি এবং কেবল যদি
$n>m$ হয়। $L$ ও $I$-এর সংযোগকে $K=L\cup I$ বললে, সেটি
\emph{চেয়ে ছোট বা সমান হওয়ার} সম্পর্ক: $Knm$ হবে যদি এবং কেবল যদি
$n \le m$ হয়। একইভাবে, $H=G \cup I$ হল \emph{চেয়ে বড় বা সমান
হওয়ার সম্পর্ক}। এই $L$, $G$, $K$ ও $H$ বিশেষ ধরনের সম্পর্ক,
যাদের \emph{ক্রম} বলে। $L$ ও $G$-এর একটি ধর্ম হল: কোনও সংখ্যাই
নিজের সঙ্গে $L$ বা $G$ সম্পর্কে থাকে না (অর্থাৎ, সমস্ত $n$-এর জন্য,
$Lnn$ কিংবা $Gnn$, কোনওটিই হয় না)। এই ধর্মবিশিষ্ট সম্পর্ককে
\emph{আত্মসম্পর্কহীন} বলে। সেটি একই সঙ্গে ক্রম হলে তাকে
\emph{কঠোর ক্রম} বলা হয়।
\end{ex}

\begin{explain}
ক্রম ও অভিন্নতা গুরুত্বপূর্ণ এবং স্বাভাবিক সম্পর্ক হলেও জোর দিয়ে বলা
দরকার যে, আমাদের সংজ্ঞা অনুযায়ী $A^{2}$-এর \emph{যে-কোনও} উপসেটই
$A$-এর উপর সম্পর্ক, সেটিকে যতই অস্বাভাবিক বা কৃত্রিম মনে হোক।
বিশেষত, $\emptyset$ যে-কোনও সেটের উপর সম্পর্ক। একে \emph{শূন্য
সম্পর্ক} বলে; কোনও উপাদানযুগলই এই সম্পর্কে থাকে না। আবার $A^{2}$
নিজেও $A$-এর উপর সম্পর্ক: প্রতিটি যুগলই এই সম্পর্কে থাকে।
একে \emph{সার্বিক সম্পর্ক} বলে। তবে
$E=\Setabs{\tuple{n, m}}{n>5 \text{ অথবা } m \times n \ge 34}$-এর
মতো কিছুও সম্পর্ক হিসেবে গণ্য হয়।
\end{explain}

\begin{prob}
  $\Pow{\{a, b, c\}}$ সেটের উপর $\subseteq$ সম্পর্কটির
  !!{element}sের তালিকা তৈরি করো।
\end{prob}

\end{document}
